\section{微分流形}

\noindent\textbf{1. 基本定义和例子}\quad
现在我们开始比较系统地讨论弯曲的空间，并在弯曲的空间上建立起微积分学。

当然，我们需要把研究的对象明确起来。一方面要求足够广泛，能包括我们可能会遇到的对象，主要是在物理学与几何学中的对象；另一方面又不能没有界限以至于会因为不能不处理一些“病态”的对象，而把我们已经掌握得很彻底的微积分理论丢开而另起炉灶。这两个方面分寸如何掌握最终是由物理学和几何学的需要决定的。科学的发展证明，一个比较适当的界限可以划到微分流形为止。

什么是一个微分流形？我们从一个实例开始。绘制地图时，我们没有办法把地球球面展开成一个平面而且保持每一个国家、地区“形状不变”。于是我们把地球分成许多小的区域，把各个区域都投影到某一个平面区域上，而且各个不同区域采用不同的投影方法，例如麦卡托（G.~Mercator）（不是第二章讲到对数函数时提到的那个麦卡托）在1569年提出的麦卡托投影法。但最易懂的是柱面投影法：把地球看作一个球面而沿赤道用一个圆柱包围起来，并从球心将球面上各点投影到柱面上，再把柱面打开成一个平面就得到一张世界全图（世界全图称为 atlas，此词来自希腊神话，是麦卡托首先用这个词的。分国分省地图都不能称为 atlas。讲一下这个故事对理解下文有好处）。这个投影法的好处是子午线和纬圈变成了互相垂直的直线族，而且在赤道附近，区域的“形状”保持得很好，但是例如北极被投影到无穷远处，而且高纬度的国家如俄罗斯、加拿大都大得出奇，很不“真实”了。另一种投影法例如球极投影（前面也讲过），如果从南极开始把球面向北极的切平面投影，则子午线变成了过北极的直线，纬圈仍是圆周，很像极坐标，它的好处是：在北极附近，图形很“真实”，而且经过北极的大圆（即子午线）都是直线，角度大小也不变，但是南极被映到无穷远处去了。因此地图学的实践是在不同区域采用不同的投影，而在采用了不同投影的两个区域之公共部分，则要研究其“迁移”（transition，这是一个数学名词而不是制图学名词）。事实上如果在第一种投影之下平面上的点用笛卡儿坐标 $(x_1,x_2)$ 表示，第二种投影之下得到 $(y_1,y_2)$，则同一点 $P$ 就既有 $(x_1,x_2)$ 又有 $(y_1,y_2)$，用箭头图示如下：

\begin{figure}[htbp]
\centering
\begin{tikzpicture}[bella figure,>=stealth,scale=.95]
  \node (P) at (0,1.7) {$P$};
  \node (X) at (-2,0) {$(x_1,x_2)$};
  \node (Y) at (2,0) {$(y_1,y_2)$};
  \draw[->] (P)--(X);
  \draw[->] (P)--(Y);
  \draw[->] (-1.35,.18)--(1.35,.18) node[midway,above] {$\varphi$};
  \draw[->] (1.35,-.18)--(-1.35,-.18) node[midway,below] {$\varphi^{-1}$};
\end{tikzpicture}
\end{figure}

从这个图看到，在 $(x_1,x_2)$ 与 $(y_1,y_2)$ 之间有一个映射 $\varphi$（其逆映射为 $\varphi^{-1}$），而用坐标表示成为
\begin{equation}
 y_1=\varphi_1(x_1,x_2),\qquad y_2=\varphi_2(x_1,x_2).
 \tag{1}
\end{equation}
这一对函数就称为“迁移函数”（transition functions，这又是数学名词而不是制图学的名词）。如果我们想在地球上做微积分，当然应该要求它们是光滑函数（但还不止于此）。从这里我们汲取了以下的研究弯曲空间的基本思想：

\begin{enumerate}[label=\arabic*.,leftmargin=2.8em]
\item 把研究对象分成许多小的部分，而每一部分都可以投影到欧氏空间 $\mathbb R^n$ 内，因而可以进行局部的讨论。
\item 在两个部分相交的地方，通过研究迁移函数来实现从一个部分到另一个部分的转换。
\end{enumerate}

现在就可以介绍微分流形的概念了。但是立刻就遇到一个问题，在讨论地球表面的各种投影以及在讨论更一般的曲面时，我们的对象是很清楚的：它是 $\mathbb R^3$ 的曲面，它的弯曲是相对于包含它的空间——所谓包含空间（ambient space）——而言的。但是本章一开始就指出，黎曼的基本思想是：空间本身就是弯曲的，用不着有包含空间作衬托。而其弯曲就表现在，其各点附近就距离而言并非均匀的，因而度量
\[
 ds^2=g_{ij}(x)\,dx^i dx^j
\]
随 $x$ 而异。现在我们也不能设想还有什么包含空间，也不知道有没有度量。因而我们一开始只有把研究对象 $M$ 规定是一个拓扑空间。

然后，我们把 $M$ 划分成许多小片，具体说，设有若干个开集——注意，拓扑空间中是可以定义开集 $\{U_\alpha\}$ 的，使
\[
 M=\bigcup_\alpha U_\alpha.
\]
对于每一个 $U_\alpha$，我们都假设有一个“投影” $\varphi_\alpha:U_\alpha\to\mathbb R^n$，从数学上说，$\varphi_\alpha$ 是一个同胚，因此 $\varphi_\alpha(U_\alpha)\subset\mathbb R^n$ 是一个开集，甚至是一个区域。$U_\alpha$ 称为一个区图（chart，这是一个古老的英文文字，下面的解释现在已不用了；一片水域或陆地的地图），而其集 $\{U_\alpha\}$ 称为 $M$ 的一个图册（atlas），这都是借用了制图学中的名称。$\varphi_\alpha(U_\alpha)\subset\mathbb R^n$，$\mathbb R^n$ 中是有坐标的，我们就借用这个坐标作为 $U_\alpha\subset M$ 中点的坐标。这种坐标称为局部坐标。$\varphi_\alpha$ 称为坐标映射（有时也就称 $\varphi_\alpha$ 为局部坐标）。我们不妨设 $U_\alpha$ 是连通的，若不然，我们可以把 $U_\alpha$ 的各个连通分支各当作一个 $U_\alpha$。因为 $\varphi_\alpha$ 是同胚，所以 $\varphi_\alpha(U_\alpha)\subset\mathbb R^n$ 也是连通开集——所谓区域就是连通开集。这样我们看到，通过同胚 $\varphi_\alpha$ 可以把 $\mathbb R^n$ 的拓扑性质——连通性、局部紧性——都移植到 $M$ 上。结果，所有我们将要定义的微分流形 $M$ 都是局部连通与局部紧的。可见，微分流形将不会是一个过于空泛的概念。

如果 $U_1\cap U_2\ne\varnothing$，则不论是否 $\varphi_1(U_1)\cap\varphi_2(U_2)\ne\varnothing$，在 $\varphi_1(U_1)$ 和 $\varphi_2(U_2)$ 的局部坐标 $(x_1,\ldots,x_n)$ 与 $(y_1,\ldots,y_n)$ 之间都会得到迁移函数作为一个映射：
\[
 \varphi_2\circ\varphi_1^{-1}:\varphi_1(U_1)\longrightarrow\varphi_2(U_2).
\]
用局部坐标来表示就是
\begin{equation}
 y_i=(\varphi_2\circ\varphi_1^{-1})_i(x_1,\ldots,x_n)
 =\psi_i(x_1,\ldots,x_n).
 \tag{2}
\end{equation}
为了使得在 $M$ 上有微分学，就需要一个至关重要的相容性条件：(2)是微分同胚。重要的是，我们不妨取 $\varphi_1(U_1),\varphi_2(U_2)$ 充分小，使得这个微分同胚条件可以用坐标来表示为：$\psi\in C^\infty,\psi^{-1}\in C^\infty$，且
\begin{equation}
 \left.\det\frac{\partial(y_1,\ldots,y_n)}{\partial(x_1,\ldots,x_n)}\right|_{\varphi_1(U_1)}\ne0.
 \tag{3}
\end{equation}
“充分小”这样的限制词是必要的，如若不然，仅由雅可比行列式(3)不为0尚不足以保证 $\varphi_2\circ\varphi_1^{-1}=\psi$ 为一微分同胚。

最后我们还要加上两个条件：首先 $M$ 是豪斯多夫空间，这不仅使我们避免了与许多“病态”的对象打交道，而且有了它，再加上第二个条件即 $M$ 适合“第二可数性公理”，即存在可数多个开集 $\{U_i;i=1,2,\ldots\}$，使 $M$ 中任一开集必是若干个（可能是可数多个）$U_i$ 之并，就可以构造出许多极为有用的重要的对象。首先是单位分解（见第三章），有了它，才能把有关 $M$ 的问题例如积分问题归结为在各个区图 $U_\alpha$ 上的相应问题。至今我们没有说，$M$ 上是否可以按照黎曼的要求构造出一个度量 $ds^2=g_{ij}(x)dx^i dx^j$，而在有了单位分解以后，就可以证明 $M$ 上确可定义黎曼的度量，证明见第5节。这以后 $M$ 之每个连通分支均各将成为一个度量空间，这一点我们则不加证明了。

有了这些讨论，我们就可以正式地给出微分流形的定义。

\noindent\textbf{定义1}\quad
设 $M$ 为一适合第二可数性公理的豪斯多夫拓扑空间，而且有由连通开集 $\{U_\alpha\}$ 所成的覆盖，使得对每个 $U_\alpha$ 存在同胚映射 $\varphi_\alpha:U_\alpha\to\mathbb R^n$，而且当 $U_\alpha\cap U_\beta\ne\varnothing$ 时迁移函数
\[
 \varphi_\beta\circ\varphi_\alpha^{-1}:\varphi_\alpha(U_\alpha)\longrightarrow\varphi_\beta(U_\beta)
\]
为一 $C^\infty$ 微分同胚，这时称 $M$ 为一 $n$ 维（$C^\infty$）微分流形，$U_\alpha$ 称为一个区图（或坐标邻域），$\{U_\alpha\}$ 称为图册，$\varphi_\alpha(U_\alpha)\subset\mathbb R^n$ 在 $\mathbb R^n$ 中的坐标称为 $M$ 中 $U_\alpha$ 之局部坐标，$\varphi_\alpha$ 称为坐标映射，$n$ 称为 $M$ 之维数。

\noindent\textbf{注1}\quad
定义中只要求 $\{U_\alpha\}$ 能覆盖 $M$，而有时需要把一切适合 $\varphi(V)\subset\mathbb R^n$ 的同胚 $\varphi$ 与连通开集 $V$ 也都列入 $\{\varphi_\alpha\}$ 和 $\{U_\alpha\}$ 之中，但要 $\varphi\circ\varphi_\alpha^{-1}$ 为微分同胚，以得到一个极大图册（maximal atlas）。极大图册的存在性是需要而且可以证明的。这种极大图册称为一个微分构造。于是产生一个问题：同一个 $M$ 上是否有唯一的微分构造。20世纪60年代，J.~Milnor 证明了 $S^7$ 上就有互不相容的微分构造存在，即一个构造中的 $\varphi_\alpha$ 与另一构造中的 $\varphi_\beta$ 能使 $\varphi_\alpha\circ\varphi_\beta^{-1}$ 不是微分同胚。也有人证明了确有没有微分构造的流形存在。20世纪80年代，英国数学家 S.~K.~Donaldson 证明了 $\mathbb R^4$ 中即有互不相容的微分构造，这件事与理论物理有密切的关系。

\noindent\textbf{注2}\quad
上面我们讲的是 $C^\infty$ 微分流形，即指迁移函数是 $C^\infty$ 函数。这是最常见的情况，一般数学文献中凡讲到光滑，总是指的 $C^\infty$，所以也时常略去 $C^\infty$（甚至光滑）字样。但实际上还有迁移函数只是 $C^k$ 函数的情况。这时流形就称为 $C^k$ 流形，特别是 $C^0$ 流形，其实并没有可微性，它的性质与其它 $C^k$ 流形可以很不相同。例如可以问，每个微分构造是否只有一个维数？确实如此，但对于 $C^k$ 流形，当 $k>0$ 时证明并不难，$k=0$ 时也是可以证明的，但要证明维数不变需证明著名的区域不变性定理：

\noindent\textbf{定理1}\quad
若 $n\ne m$，则 $\mathbb R^n$ 之一开集不能同胚于 $\mathbb R^m$ 的开集。

这个著名定理是布劳威尔证明的，要用到很细致的同调理论工具，但在 $k>0$ 时却不难证明。由此，我们至少看到，维数是刻画一个流形的整体不变量；对一个微分流形找出尽可能多的整体不变量是很重要的事。这类问题我们就说到这里为止。

有时，迁移函数是实解析函数（一个函数在某点附近是实解析的，即在此点附近可以展开为收敛的幂级数），实解析函数类时常记为 $C^\omega$，实解析流形也就记为 $C^\omega$ 流形。读者会问，如果把上面讲的 $n$ 维实空间 $\mathbb R^n$ 换成 $n$ 维复空间，而迁移函数从可微函数换成全纯函数，是否可以得到相应的复流形理论？是，而且复流形与实流形很不相同，我们完全不能涉及。

$M$ 上有了微分结构以后，就可以在其上建立种种有关微分学的概念。这里一定要注意，所有这些概念都必须是与局部坐标之选取无关的。我们先确定一下什么是 $M$ 上的可微函数 $f$。如果 $f$ 定义在 $M$ 上某点 $P$ 之邻域中，不失一般性，可以设此邻域即为一坐标邻域 $U$，于是 $f\circ\varphi^{-1}$ 即成为 $\varphi(U)\subset\mathbb R^n$ 上的函数。如果 $f\circ\varphi^{-1}$ 是 $C^\infty$ 可微的，就说 $f$ 在 $P$ 附近是 $C^\infty$ 可微的。这个定义显然与 $U$ 和 $\varphi$ 的选择无关，因为若选 $P$ 的另一个邻域 $V$ 与相应的局部坐标 $\psi$，但在 $U\cap V$ 上
\[
 \psi\circ\varphi^{-1}:\varphi(U\cap V)\longrightarrow\psi(U\cap V)\subset\psi(U)
\]
是一个微分同胚，因而是 $C^\infty$ 的，$f\circ\varphi^{-1}$ 也是 $C^\infty$ 函数，所以
\[
 f\circ\psi^{-1}=(f\circ\varphi^{-1})\circ(\varphi\circ\psi^{-1})
\]
至少在 $P$ 的一个较小的邻域 $U\cap V$ 上是 $C^\infty$ 可微的。

下面我们来举一些 $C^\infty$ 微分流形的例子。

\noindent\textbf{例1}\quad
$M=\mathbb R^n$。这是平凡的，因为现在的图册 $\{U_\alpha\}$ 只含一个元 $U=\mathbb R^n$ 即可，而坐标映射就是恒等映射：$\varphi_\alpha=\mathrm{id}$。我们这里没有去讨论什么是极大图册——微分构造，讨论那个问题对我们目前并无好处。

其次令 $M$ 为 $\mathbb R^n$ 中任一开集 $U$，这时仍有 $\{U_\alpha\}=\{U\}$，$\varphi_\alpha=\mathrm{id}|_U$，我们说 $U$ 是 $\mathbb R^n$ 的开子流形。这个概念与下面要讲到的闭子流形是不一样的，开子流形的维数仍为 $n$。

读者会问，半空间 $H^n=\{x\in\mathbb R^n;x_n\ge0\}$ 是不是微分流形？一方面我们很需要把它纳入我们的视野，因为例如在积分理论中考虑斯托克斯（Stokes）定理（格林公式是它的特例）时，必须要讨论 $H^n$ 上的积分。另一方面它又不是流形，因为它有一个“边缘” $x_n=0$，其上的点没有任何可以同胚于 $\mathbb R^n$ 的邻域。因此，我们要以 $H^n$ 为模型来扩大微分流形的定义。

\noindent\textbf{定义2}\quad
设 $M^+$ 是一适合第二可数性公理的豪斯多夫拓扑空间，$\{U_\alpha\}$ 是覆盖 $M^+$ 的开集族，$\bigcup U_\alpha=M^+$，而对每个 $U_\alpha$ 均有同胚 $\varphi_\alpha:U_\alpha\to\mathbb R^n$ 或 $H^n$，使得当 $U_\alpha\cap U_\beta\ne\varnothing$ 时，$\varphi_\beta\circ\varphi_\alpha^{-1}:\varphi_\alpha(U_\alpha)\to\varphi_\beta(U_\beta)$ 为一微分同胚，这时称 $M^+$ 为带边的微分流形。$\{x\in H^n;x_n=0\}$ 在 $M^+$ 中的原像称为 $M^+$ 的边缘，记作 $\partial M^+$。

\noindent\textbf{注}\quad
一个带边流形 $M^+$ 之边缘与一个拓扑空间的子空间的边界点是不同的概念。例如，$\mathbb R^2$ 中的去了圆心的闭圆盘 $0<r\le1$ 的边缘只是单位圆周 $r=1$，但其边界点之集合还多了一个点 $r=0$——圆心。

注意，这个定义中 $M^+$ 就是拓扑空间的全空间，因此不存在在 $M^+$ 以外的点的问题。也因此，$U_\alpha$ 作为 $M^+$ 中的开集允许含有 $\partial M^+$ 上之点。但是 $H^n$ 又是 $\mathbb R^n$ 的子空间，并赋有子空间拓扑，因此 $\{x_n=0\}$ 上之点虽然就 $H^n$ 是 $\mathbb R^n$ 之子集而言是 $H^n$ 之边界点，按子空间拓扑而言则可以含于 $H^n$ 之开集 $\varphi_\alpha(U_\alpha)$ 中。由同样的理由，$\varphi_\beta\circ\varphi_\alpha^{-1}$ 可以一直光滑到 $x_n=0$ 处，我们通常就说，$\varphi_\beta\circ\varphi_\alpha^{-1}$ “光滑到边”。现在我们问，把 $H^n$ 作为 $\mathbb R^n$ 之子集来看，可否找到“更大”的子集 $K$，包含 $H^n$ 而且跨越到 $x_n<0$ 处，使 $\varphi_\beta\circ\varphi_\alpha^{-1}$ 可以光滑地拓展到 $K$ 上？一般说来，设 $f(x)$ 是闭集 $H^n$ 上的光滑函数，$K\supset H$ 是一开集，能否找到一个定义在 $K$ 上的光滑函数 $F(x)$，使在 $H$ 上 $F(x)=f(x)$？$F(x)$ 称为 $f(x)$ 由 $H$ 到 $K$ 上的光滑延拓。是否可能延拓与 $\partial H$ 的构造有密切的关系。作为带边流形的模型 $H^n$ 之边缘构造特别简单，就只是一个超平面 $x_n=0$，比较一般的情况，则可以考虑 $M^+$ 局部也可以用 $f(x)\ge0$ 来刻画，这里 $f(x)\in C^\infty$，而且 $\mathrm{grad}\,f\ne0$，且 $M^+$ 局部地位于 $\partial M^+=\{x;f(x)=0\}$ 之一侧。在这个情况下有著名的 Seeley 延拓定理：

\noindent\textbf{Seeley 延拓定理}\quad
设 $f(x)\in D(H^n)$，则必有一个连续映射 $\sigma:D(H^n)\to D(\mathbb R^n)$，使
\[
 (\sigma f)(x)=f(x),\qquad x\in H^n.
\]
由于有这个定理，上述迁移函数 $\varphi_\beta\circ\varphi_\alpha^{-1}$ 可以在 $x_n<0$ 的某一区域中也光滑，且是微分同胚，这对我们讨论带边流形有很大的方便。

这个定理的证明我们就略去了，只是要提醒一下，如果 $\partial M^+$ 不是如此简单，哪怕 $M^+$ 例如是一个象限，角点的出现也会使函数的延拓困难不少。

\noindent\textbf{例2\quad 无奇点的光滑超曲面}\quad
设有 $\mathbb R^n$ 中一个光滑的超曲面，它由一个方程
\begin{equation}
 f(x_1,\ldots,x_n)=0
 \tag{4}
\end{equation}
定义，这里 $f$ 是其自变量在 $\mathbb R^n$ 之某区域 $\Omega$ 内的 $C^\infty$ 函数，我们设它是无奇点的，即在 $\Omega$ 上
\[
 \mathrm{grad}\,f(x)\ne0,
\]
因此在每一点至少有一个偏导数不为0。设在某个开集 $U$ 中，$\partial f/\partial x_n\ne0$ 而可以应用隐函数定理得出
\[
 x_n=\varphi(x_1,\ldots,x_{n-1}),\qquad (x_1,\ldots,x_{n-1})\in\omega.
\]
这里，适合以上方程的 $(x_1,\ldots,x_{n-1},x_n)\in U$，从而所有这种 $U$ 构成超曲面(4)的开覆盖，而且有坐标映射
\[
 \Phi:\{(x_1,\ldots,x_{n-1},\varphi(x_1,\ldots,x_{n-1}))\in U,
 (x_1,\ldots,x_{n-1})\in\omega\}\longrightarrow\omega\subset\mathbb R^{n-1}.
\]
于是我们需要证明的只是，例如有两个这样的 $U$ 和 $V$，在 $V$ 中 $\partial f/\partial x_{n-1}\ne0$，且有相应的坐标映射
\[
 \begin{aligned}
 \Psi:\{&(x_1,\ldots,x_{n-2},\psi(x_1,\ldots,x_{n-2},x_n),x_n),\;\\
 &(x_1,\ldots,x_{n-2},x_n)\in\Omega\}\longrightarrow\Omega\subset\mathbb R^{n-1}.
\end{aligned}
\]
则 $\Psi\circ\Phi^{-1}:\omega\to\Omega$ 是迁移函数。\linebreak[4] 其自变量是 $(x_1,\ldots,x_{n-1})$，而函数值是 $(x_1,\ldots,x_{n-2},x_n)$。但是很清楚，因为
\[
 \Psi\circ\Phi^{-1}:(x_1,\ldots,x_{n-1})\longmapsto(x_1,\ldots,x_{n-2},x_n)
\]
可以写为
\[
 x_j=x_j,\quad j=1,2,\ldots,n-2,
 \qquad x_n=\varphi(x_1,\ldots,x_{n-1}),
\]
且
\[
 \frac{\partial x_n}{\partial x_{n-1}}
 =-\frac{\partial f/\partial x_{n-1}}{\partial f/\partial x_n}\ne0,
\]
所以它当 $\omega,\Omega$ 充分小时是微分同胚。总之上述超曲面是一个 $n-1$ 维微分流形。

将它推广，在 $\mathbb R^n$ 中由 $m$（$m<n$）个方程
\begin{equation}
 f_j(x_1,\ldots,x_n)=0,\qquad j=1,\ldots,m
 \tag{5}
\end{equation}
定义的几何轨迹是一个 $n-m$ 维微分流形，这里 $f_j$ 是 $C^\infty$ 函数，而且 $m\times n$ 矩阵
\[
 \begin{pmatrix}
 \dfrac{\partial f_1}{\partial x_1}&\cdots&\dfrac{\partial f_1}{\partial x_n}\\
 \vdots&&\vdots\\
 \dfrac{\partial f_m}{\partial x_1}&\cdots&\dfrac{\partial f_m}{\partial x_n}
 \end{pmatrix}
\]
之秩为 $m$。

这个例子告诉我们，$\mathbb R^n$ 中的曲面(5)是微分流形。反过来，也有定理说明任一个 $k$ 维微分流形都可以嵌入在充分高维的欧氏空间 $\mathbb R^{2k+1}$ 中。这就是说，为了研究微分流形，只需讨论本例就够了，这时，$\mathbb R^n$ 成了此流形的包含空间。但是我们并不采取这种观点，因为它使得许多很深刻的思想得不到充分表现，而且，维数很高的空间带来的几何复杂性不见得比不用包含空间来抽象地讨论微分流形更少。

\noindent\textbf{例3\quad 黎曼球面和射影空间}\quad
上一章中我们在 $z$ 平面 $\mathbb C$（即 $(x,y)$ 平面 $\mathbb R^2$）上引进了一个无穷远点，并通过球极射影把它变成一个球面 $S^2$，称为黎曼球面。从球极射影的公式可以看出这个射影其实是可微的。现在我们要证明黎曼球面是一个2维微分流形，而球极射影其实就是坐标映射。

\begin{figure}[htbp]
\centering
\begin{tikzpicture}[bella figure,>=stealth,scale=.82]
  \draw (0,0) circle (2.15);
  \draw[dashed] (-2.1,0)--(2.1,0);
  \draw (-3.5,-.35)--(3.5,-.35)--(4.3,.55)--(-2.7,.55)--cycle;
  \coordinate (N) at (0,2.15);
  \coordinate (S) at (0,-2.15);
  \coordinate (P) at (1.35,.95);
  \coordinate (Q) at (2.95,-.35);
  \fill (N) circle (1.2pt) node[above] {$N$};
  \fill (S) circle (1.2pt) node[below] {$S$};
  \fill (P) circle (1.2pt) node[right] {$(\xi,\eta,\zeta)$};
  \fill (Q) circle (1.2pt) node[below right] {$(x,y)$};
  \fill (0,0) circle (1.1pt) node[left] {$O$};
  \draw (N)--(Q);
\end{tikzpicture}
\caption*{图7-2-1}
\end{figure}

由图7-2-1可以看出，$S^2=U_1\cup U_2$，$U_1=S^2\setminus\{N\}$，$N$ 是北极 $(0,0,1)$；$U_2=S^2\setminus\{S\}$，$S$ 是南极 $(0,0,-1)$。在 $U_1$ 中作球极射影 $\varphi_1:U_1\to\mathbb R^2$，$\mathbb R^2$ 是过球心的赤道平面，并把一点 $P(\xi,\eta,\zeta),\zeta\ne1$ 映到 $\mathbb R^2$ 上的 $(x,y)$，很容易看到
\[
 x=\frac{\xi}{1-\zeta},\qquad y=\frac{\eta}{1-\zeta},\qquad \zeta\ne1.
\]
这就是 $U_1\to\mathbb R^2$ 的坐标映射。同样，若从南极 $S$ 开始作球极射影，则点 $(\xi,\eta,\zeta)$ 对应的点 $(x_1,y_1)$ 应为
\[
 x_1=\frac{\xi}{1+\zeta},\qquad y_1=\frac{\eta}{1+\zeta},\qquad \zeta\ne-1.
\]
这就是 $U_2\to\mathbb R^2$ 的坐标映射。由此就可以得到，除在 $(x,y)=(0,0)$ 处以外，
\[
 x_1=\frac{x}{x^2+y^2},\qquad y_1=\frac{y}{x^2+y^2}.
\]
而在 $U_1\cap U_2$ 中恰好没有南北极，因此不会有 $(x,y)=(0,0)$ 和 $(x_1,y_1)=(0,0)$，所以上式就是 $U_1\cap U_2$ 中的迁移函数。由此我们看到，黎曼球面是一个微分流形。其实，黎曼球面只不过是普通的复平面即 $(x,y)$ 平面 $\mathbb R^2$ 添加一个无穷远点而已，但是它却成了一个与 $S^2$ 同胚的球面，与 $\mathbb R^2$ 的性质大相径庭：前者是紧的，而后者只有局部紧，尽管局部地看来都是通常的平面。如果我们进一步再看射影平面 $RP^2$，还会感受到更多的区别。

射影平面其实也只是普通平面再添加一些理想的元素。这一次是添加了一条无穷远直线却又使它们的性质与黎曼球面大不相同了。然而 $RP^2$ 仍然只是一个2维微分流形。

从上一章看到，定义 $RP^2$ 有多种方法，其一是在 $\mathbb R^3$ 中把过原点的直线当成一个元素，每一条这样的直线可以用其上非原点的一个点 $(x_1,x_2,x_3)$ 来定义，这里 $x_1,x_2,x_3$ 不能同时为0。不过，若 $(x_1,x_2,x_3)$ 与 $(x'_1,x'_2,x'_3)$ 只相差一个常数因子，即
\[
 x_i=\lambda x'_i,\quad \lambda\ne0,
 \qquad\text{或}\qquad
 x'_i=\mu x_i,\quad \mu\ne0,
\]
这两点要算成等价的：$(x_1,x_2,x_3)\sim(\lambda x_1,\lambda x_2,\lambda x_3)$。而这些直线的集合即 $\mathbb R^3$ 中这些点的等价类的集合，我们就称 $\mathbb R^3/\sim$ 为2维的实射影平面：$RP^2=\mathbb R^3/\sim$。上一章里我们在 $RP^2$ 上这样来定义一个拓扑，使 $RP^2$ 与球面 $S^2$（经过一些改造）同胚。现在进一步问，能否用类似的方法在 $RP^2$ 上引入微分结构，使得在其上可以建立起微分学？

上一章里，我们是令通过原点的直线与 $S^2$ 相交，而将直线之集合映到 $S^2$ 上。但每一条这样的直线与 $S^2$ 有两个交点且为对径点，因此要把 $S^2$ 上的对径点视为一点，这样得到射影平面的几何形象，也就是上文说的 $S^2$“经过一点改造”的意思。这样作会发现，$RP^2$ 与一些很奇特的几何图形例如莫比乌斯带有关系。也可以这样去讨论 $RP^2$ 的微分构造，但这需要把商拓扑的概念和方法讲得更明确，因此我们采用另一个方法。

讲解射影平面最简单的方法是引入 $\mathbb R^2$ 上的所谓“齐次坐标”。我们先引入 $\mathbb R^3$，即在 $\mathbb R^2$ 上方多加一个 $z$ 轴，于是 $\mathbb R^2$ 就成了 $z=0$。我们先把这个平面向上抬起成为 $z=1$，于是原来的 $(x,y)$ 点变成 $(x,y,1)$。把它与原点 $(0,0,0)$ 连成直线，这样一来，$\mathbb R^2$ 与 $\mathbb R^3$ 中过原点的直线族对应起来。上面已经说过这个直线可以用 $(x,y,z)$（$x,y,z$ 不同时为0）之等价类 $(x,y,z)\sim(\lambda x,\lambda y,\lambda z)$，$\lambda\ne0$ 对应起来。于是 $(x,y)$ 点也就与 $(x,y,1)$ 之等价类对应起来了。我们选此等价类之任一元 $(x_1,x_2,x_3)$ 作为代表元，但要求 $x_1,x_2,x_3$ 不同时为0，例如 $(x,y)$ 对应于 $(x,y,1)$ 的等价类，而 $x_1=x,x_2=y,x_3=1$，于是 $x=x_1/x_3,y=x_2/x_3$。正如复平面变成黎曼球面是通过附加无穷远点（即北极）而得，我们现在也对 $\mathbb R^2$ 添加“许多”点成为一个新的集合，即加上 $\{(x_1,x_2,0)\}$，并把这个新集合称为射影平面 $RP^2$。它的点是坐标 $(x_1,x_2,x_3)$（坐标不同时为0），称为射影平面中一点的齐次坐标。当 $x_3\ne0$ 时，由它们可以得出普通的笛卡儿坐标 $x=x_1/x_3,y=x_2/x_3$。$RP^2$ 中每个点并不是恰有一个齐次坐标，而是有齐次坐标的恰好一个等价类。齐次坐标方法很早在数学中就通行了。例如当我们研究一个多项式 $P(x,y)=0$ 时，就可以引入齐次坐标，并考虑一个三个变元的齐次多项式
\[
 x_3^m P\left(\frac{x_1}{x_3},\frac{x_2}{x_3}\right)=F(x_1,x_2,x_3)=0,
\]
$m$ 是 $P$ 的次数。齐次多项式的好处不仅在于形式整齐，而且它允许添加一些理想的元素。例如直线 $ax+by+c=0$ 现在变成 $ax_1+bx_2+cx_3=0$。这样一来就增加了一个研究对象 $x_3=0$，它相应于 $a=b=0,c\ne0$。本来我们是不应该考虑 $x_3=0$ 的，因为 $a=b=0$ 而 $c\ne0$ 与 $ax+by+c=0$ 是矛盾的，上面用 $x_3^m$ 去乘 $P(x_1/x_3,x_2/x_3)$ 也是不合法的。但是若我们允许这样做了就相当于把集合 $\{(x_1,x_2,0)\}$ 即 $x_3=0$ 也纳入我们的讨论之中，就相当于把 $\mathbb R^2$ 变成 $RP^2$。这种做法的合理性在于：$x_3=0$ 和 $ax_1+bx_2+cx_3=0$ 一样是一个一次方程，而一次方程代表一条直线，所以我们其实是在 $\mathbb R^2$ 中添上了一条直线——无穷远直线。这条直线经过一个适当的射影变换即齐次坐标的非奇异线性变换，可以化为任一指定的普通的直线，所以从射影几何的角度来看 $x_3=0$ 与其它的直线并无本质的不同。

现在我们可以来证明 $RP^2$ 确实是一个微分流形。事实上
\begin{equation}
 RP^2=\bigcup_{i=1}^3U_i
 =\bigcup_{i=1}^3\{(x_1,x_2,x_3);x_i\ne0\}.
 \tag{6}
\end{equation}
$U_i=\{(x_1,x_2,x_3);x_i\ne0\}$ 显然是一个开集。从 $U_i$ 到 $\mathbb R^2$ 之同胚定义为
\[
 \varphi_i:U_i\longrightarrow\mathbb R^2,
 \quad (x_1,x_2,x_3)\longmapsto
 \left(\frac{x_1}{x_i},\ldots,\frac{x_{i-1}}{x_i},
 \frac{x_{i+1}}{x_i},\ldots,\frac{x_3}{x_i}\right),
\]
当 $i=1$ 或3时，$x_{i-1}$ 或 $x_{i+1}$ 无意义，这时要稍作些修改，总之使分子上不再出现 $x_i$。余下的只要证明 $\varphi_i\circ\varphi_j^{-1}$ 在 $\varphi_j(U_i\cap U_j)$ 上是微分同胚即可。不妨看 $\varphi_2\circ\varphi_1^{-1}$\footnote{校勘：此处底本写作 $\varphi_1\circ\varphi_2^{-1}$，但以下以 $z=\varphi_1(\cdot)$、$w=\varphi_2(\cdot)$ 表示两组坐标，并给出 $w$ 关于 $z$ 的公式，因此这里应为 $\varphi_2\circ\varphi_1^{-1}$。}。在 $U_1\cap U_2(\ne\varnothing)$ 上，
\[
 \varphi_1(U_1\cap U_2)
 =\left\{\left(\frac{x_2}{x_1},\frac{x_3}{x_1}\right);x_1x_2\ne0\right\}
 =\{(z_1,z_2)\},
\]
而
\[
 \varphi_2(U_1\cap U_2)
 =\left\{\left(\frac{x_1}{x_2},\frac{x_3}{x_2}\right);x_1x_2\ne0\right\}
 =\{(w_1,w_2)\},
\]
于是 $\varphi_2\circ\varphi_1^{-1}$ 可以写为
\[
 w_1=\frac1{z_1},\qquad w_2=\frac{z_2}{z_1}.
\]
但因在 $U_1\cap U_2$ 中 $z_1=x_2/x_1\ne0$，所以上式很明显是微分同胚。于是 $RP^2$ 是一个微分流形。

比较一下黎曼球面与 $RP^2$，就知道，局部地看，它们是非常相似的，都是 $\mathbb R^2$。其实几何微分流形必为局部欧氏空间。但是整体来看却极不相同；$RP^2$ 虽然与黎曼球面同样为紧的、连通空间，却与后者不同，它是不可定向的，一般地说来，$RP^n$ 当 $n$ 为偶数时必为不可定向的，而当 $n$ 为奇数时却是可以定向的。那么 $RP^2$ 整体看起来像什么样子？这就很不好说，因为它只能嵌入（这两个字也要解释，见后文）在 $\mathbb R^4$ 中。$\mathbb R^5$ 是一个什么样的空间就不是人们的直觉能够想得到的了。因此，整体地研究微分流形我们得一整套其它方法。这就是拓扑学的主题了，它将应用同调、同伦等等一整套方法，而我们这一章只能介绍一点微分流形上的微积分学罢了。

\noindent\textbf{例4\quad 黎曼曲面}\quad
本书中我们讲了许多不同的微分流形，不止于弯曲的曲面。下面我们再讲两种在本书中涉及的微分流形，其一是黎曼曲面。在第二章中我们就复变量的多值函数 $w=\sqrt z$ 介绍了黎曼曲面的概念。它是以我们的直觉为基础的，因而很难解释它如何穿过其自身。现在我们从微分流形的角度来看 $w=\sqrt z$ 的黎曼曲面是如何构造出来的。首先，我们把 $z$ 平面看成是 $\mathbb R^2$ 即 $(x,y)$ 平面。在其上取一点 $z_0\ne0$，并作以 $z_0$ 为中心圆周过原点的圆盘 $U_0$，它是一个开集。在这个开集中考查 $w=\sqrt z$。$\sqrt{z_0}$ 有两个可能的值，相差一个符号，任取其中之一，记作 $(z_0)^{1/2}$，则有以下的二项级数展开式（见第三章）：
\begin{equation}
 \begin{aligned}
 w=\sqrt z
 &= (z_0)^{1/2}\left(1+\frac{z-z_0}{z_0}\right)^{1/2}\\
 &= (z_0)^{1/2}\sum_{n=0}^{\infty}A_n\xi_0^n,
 \qquad \xi_0=\frac{z-z_0}{z_0}.
 \end{aligned}
 \tag{7}
\end{equation}
它在 $|\xi_0|<1$ 即当 $z\in U_0$ 时收敛。因此(7)是 $w=\sqrt z$ 在 $U_0$ 中适用的表达式，或者说(7)是代数方程 $w^2-z=0$ 的一个解，另一解则是 $w=-(z_0)^{1/2}\sum_{n=0}^\infty A_n\xi_0^n$。我们在 $U_0$ 中任取一点 $z_1$，并用(7)式算出并记作 $(z_1)^{1/2}$，于是又有
\begin{equation}
 w=(z_1)^{1/2}\sqrt{1+\frac{z-z_1}{z_1}}
 =(z_1)^{1/2}\sum_{n=0}^{\infty}A_n\xi_1^n,
 \qquad \xi_1=\frac{z-z_1}{z_1}.
 \tag{8}
\end{equation}
这个表达式适合于圆盘 $U_1$（以 $z_1$ 为心、边界过原点）中（图7-2-2），而且由于 $(z_1)^{1/2}$ 的选取方法知，在 $U_0\cap U_1$ 中，(7)与(8)定义相同的函数。第三章中指出，(8)称为(7)在 $U_1$ 中的解析延拓，于是由(7)与(8)共同定义的函数之定义域是 $U_0\cup U_1$。仿此进行下去，可能在作出 $U_{N-1}$ 时就发现 $z_0\in U_{N-1}$。于是我们可以选 $z_N$ 和 $U_N$ 即为 $z_0,U_0$。但这时会得到 $(z_N)^{1/2}=-(z_0)^{1/2}$，这一点在第三章中就已讲过。所以，如果我们真的令 $U_N=U_0$，则在其上会定义两个函数，一是(7)式，另一个是(7)的反号，因此无法在 $\bigcup_{i=0}^{N}U_i$ 上定义一个单值的解析函数 $w=\sqrt z$。为此，我们只好取另一个区域亦为 $U_N$，尽管它作为 $\mathbb R^2$ 的区域仍与 $U_0$ 相同，而且 $z_N=z_0$，却是另一个区域，并且从 $U_N$ 开始再用级数
\begin{equation}
 w=(z_N)^{1/2}\sum_{n=0}^{\infty}A_n\xi_N^n,
 \qquad \xi_N=\frac{z-z_N}{z_N}
 \tag{9}
\end{equation}
作为 $w=\sqrt z$ 在 $U_N$ 中的定义（注意(7)与(9)的系数 $A_n$ 是相同的二项系数）。如果这样作下去，则在作出了 $U_{2N-1}$ 时，就会发现 $z_{2N}=z_N=z_0$，但这时 $(z_{2N})^{1/2}=-(z_N)^{1/2}=(z_0)^{1/2}$，于是我们又回到了出发点。这样做下去，可能得到许多 $U_i$，而令 $M=\bigcup U_i$，则 $w=\sqrt z$ 是定义在 $M$ 上的单值函数。在 $M$ 中可能有某些 $U_i$ 与 $U_j$ 连续条件 $U_i\cap U_j\ne\varnothing$，例如 $U_{2N-1}\cap U_0\ne\varnothing$，$U_0\cap U_1=U_{2N}\cap U_{2N+1}\ne\varnothing$，但是 $U_0\cap U_N=\varnothing$，这样的做法自然与微分流形的定义十分相近了，只差迁移函数还没有做出来。

\begin{figure}[htbp]
\centering
\begin{tikzpicture}[bella figure,scale=.75]
  \draw[->] (-4.8,0)--(4.8,0);
  \draw[->] (0,-3.2)--(0,3.5);
  \fill (0,0) circle (1.2pt) node[below left] {$0$};
  \foreach \x/\y/\r/\lab in {-2.7/1.1/1.55/$U_0$,-1.2/1.9/1.55/$U_1$,0.4/1.4/1.55/$U_2$,1.7/.4/1.55/$U_i$,2.6/-1.0/1.55/$U_{N-1}$,-.7/-.9/1.55/$U_N$}{
    \draw (\x,\y) circle (\r);
    \node at (\x+.8,\y+.9) {\lab};
  }
  \fill (-2.2,1.0) circle (1.2pt) node[above left] {$z_0$};
  \fill (-.8,1.7) circle (1.2pt) node[above] {$z_1$};
  \fill (1.2,.9) circle (1.2pt) node[right] {$z_i$};
  \fill (2.2,-.6) circle (1.2pt) node[right] {$z_{N-1}$};
\end{tikzpicture}
\caption*{图7-2-2}
\end{figure}

如果 $U_i\cap U_j\ne\varnothing$，则在 $U_i\cap U_j$ 中因为局部坐标 $\xi_i$ 与 $\xi_j$ 适合
\[
 \xi_i=\frac{z-z_i}{z_i},\qquad \xi_j=\frac{z-z_j}{z_j},
 \qquad z_i\ne0,\ z_j\ne0,
\]
所以
\begin{equation}
 \xi_j=\frac{z_i}{z_j}(1+\xi_i)-1,
 \qquad
 \xi_i=\frac{z_j}{z_i}(1+\xi_j)-1.
 \tag{10}
\end{equation}
这就是我们需要的迁移函数。分开 $\xi_i$ 与 $\xi_j$ 的实虚部，容易看到，上式是由 $\mathbb R^2$ 到 $\mathbb R^2$ 的微分同胚。但是更值得注意的是，作为复变量 $\xi_i$ 与 $\xi_j$ 的关系来看，(10)表示 $\xi_i$ 是 $\xi_j$ 的全纯函数，其逆亦然。仿照这样的方法我们可以研究代数方程
\[
 f(z,w)=0
\]
（所谓代数方程即指 $f$ 是 $z$ 与 $w$ 的多项式）的解 $w=w(z)$。这些解一般地是 $z$ 的多值函数。仿照这个特例，而可以定义一般的黎曼曲面如下：

\noindent\textbf{定义3}\quad
设 $M$ 为一个2维微分流形，$\{(U_\alpha,\varphi_\alpha)\}$ 是其一个图册，如果当 $U_\alpha\cap U_\beta\ne\varnothing$ 时，$\varphi_\beta\circ\varphi_\alpha^{-1}:\varphi_\alpha(U_\alpha)\to\varphi_\beta(U_\beta)$ 是黎曼球面上的区域 $\varphi_\alpha(U_\alpha)$ 中的全纯函数，则称 $M$ 为一黎曼曲面。

至此，我们看到一个全纯函数可以定义在一个微分流形上，只要它有黎曼曲面构造即可，而不必要定义在复平面或黎曼球面上。因此也就不必问这个黎曼曲面是怎样穿过它自身的。问这个问题其实是问这个黎曼曲面可否在 $\mathbb C$ 上“实现”。这是不可能的，因为一般说来，一个2维微分流形要 $\mathbb R^5$ 才能把它装得下——即把它嵌入在 $\mathbb R^5$ 中，由此也就可见，微分流形概念怎样扩大了我们的视野。

\noindent\textbf{例5\quad 某些矩阵类}\quad
上一节中我们介绍了一些在数学与物理学中起重要作用的变换群，特别重要的有 $O(n),SO(n),U(n),SU(n)$ 等等。当时我们指出，它们都是重要的李群。李群是既有微分构造的微分流形，又是有群构造的群，而且这两种构造应该是相容的（相容性的确切意义这里不能讲了）。现在我们来看一下这些变换群的微分流形构造。

首先，我们把这些变换都看成 $n\times n$ 矩阵。于是每一个矩阵都对应于 $\mathbb R^{n^2}$ 中的一个点，而一个特殊的变换群，将成为 $\mathbb R^{n^2}$ 中的一个流形。如我们在例2中看到的那样，这些流形将由 $\mathbb R^{n^2}$ 中的很多个方程来定义。例如 $SO(3)$ 是适合以下条件的 $3\times3$ 矩阵
\[
 A=\begin{pmatrix}
 a_{11}&a_{12}&a_{13}\\
 a_{21}&a_{22}&a_{23}\\
 a_{31}&a_{32}&a_{33}
 \end{pmatrix}
\]
定义的。这里
\begin{equation}
 \det A=1,
 \tag{11}
\end{equation}
\begin{equation}
 \sum_{k=1}^{3}a_{ik}a_{jk}=\delta_{ij},
 \qquad i,j=1,2,3.
 \tag{12}
\end{equation}
问题是，这些条件是否互相独立的？又是不是足够的？例如上面这组条件是说 $A$ 之各行成为一个 o.n. 系，那么是不是还应增加关于各列的条件
\[
 \sum_{k=1}^{3}a_{ki}a_{kj}=\delta_{ij},\qquad i,j=1,2,3
\]
才够？事实上是不必要的，有关于行的(12)就够了。但是即令在(12)式中还有多余的，(12)中有9个条件，例如 $i=1,j=2$ 以及 $i=2,j=1$ 其实是一个条件，所以这9个条件中实际上独立的最多只有6个。再有 $\det A=1$。事实上，由(12)已知 $A\in O(n)$，从而 $\det A=\pm1$，上一节已指出 $\det A=1$ 表示在变换 $A$ 之下手征不变；如果适合手征不变的正交变换构成一个流形，则改变手征的正交变换将成为另一个流形。这样，$O(n)$ 其实是两个流形，而 $\det(A)=1$ 只表示由这两个流形中选取了一个。所以条件(11)并不是与(12)互相独立的第7个条件。

从这些分析可以看到，我们现在遇到的困难将是弄清究竟有多少个独立的条件。这些问题在研究这些群时都是最起码的问题，我们在这里不去讨论它。我们只来看一个最简单的特例，即证明 $SL(n,\mathbb R)$ 是一个 $n^2-1$ 维微分流形。$SL(n,\mathbb R)$ 称为特殊线性群，即适合 $\det A=1$ 的 $n\times n$ 实矩阵所成之群，$GL(n,\mathbb R)$ 称为一般线性群，即适合 $\det A\ne0$ 的 $n\times n$ 实矩阵所成之群。这个证明其实是很简单的。

把一个矩阵 $A$ 看成 $\mathbb R^{n^2}$ 中的一个点，则 $SL(n,\mathbb R)$ 就是 $\mathbb R^{n^2}$ 中由方程 $\det A=1$ 所定义的一个子集。$\det A$ 是 $n^2$ 个自变量 $a_{ij}$ 所成的一个很特殊的多项式。要看它是否定义一个流形就要看在其上一点 $a_{ij}$ 处（这些 $a_{ij}$ 适合方程 $\det A=1$）是否没有奇点，即 $\partial(\det A)/\partial a_{ij}$ 是否不会全为0。但这是很容易弄清的。令 $a_{ij}$ 在 $A$ 中的代数余子式为 $A_{ij}$，显然其中不含 $a_{ij}$，容易算出
\[
 \frac{\partial}{\partial a_{ij}}(\det A)=0
 \quad\Longleftrightarrow\quad
 A_{ij}=0,\ \forall i,j.
\]
但若在一点处一切 $A_{ij}=0$，则在该点由拉普拉斯展开式
\[
 \det A=\sum_{j=1}^{n}a_{ij}A_{ij},\qquad i\text{ 为固定的}
\]
知 $\det A=0$，而这与 $\det A=1$ 相矛盾。

所以，$SL(n,\mathbb R)$ 是一个 $n^2-1$ 维微分流形。关于 $SO(n)$ 将在下面讨论。

现在讨论微分流形间的映射。设 $M$ 和 $N$ 是两个微分流形，其维数分别为 $m$ 和 $n$，而 $f$ 是由 $M$ 到 $N$ 中的映射，且 $f(P)=Q$。为了定义映射的可微性，我们取 $P$ 和 $Q$ 的坐标邻域 $U$ 与 $V$ 以及相应的坐标映射 $\varphi$ 和 $\psi$。于是
\[
 \psi\circ f\circ\varphi^{-1}:\varphi(U)\longrightarrow\psi(V)
\]
是由 $\mathbb R^m$ 中的一个区域 $\varphi(U)$ 到 $\mathbb R^n$ 中的区域 $\psi(V)$ 的映射，在局部坐标下就是一组 $n$ 个 $m$ 元函数。

\noindent\textbf{定义4}\quad
设 $M,N,f$ 等均如上所述，若任一点 $P\in M$ 以及 $Q=f(P)\in N$ 均有坐标邻域 $U$ 与 $V$ 以及相应坐标映射 $\varphi,\psi$ 使 $\psi\circ f\circ\varphi^{-1}:\varphi(U)\subset\mathbb R^m\to\psi(V)\subset\mathbb R^n$ 为 $C^\infty$ 光滑函数，则称 $f$ 为 $C^\infty$ 可微映射。

这个定义与局部坐标的选择无关，因为若有另一组坐标邻域与坐标映射 $(\varphi_1,U_1),(\psi_1,V_1)$，则在 $\varphi_1(U_1\cap U)$ 中，
\[
 \psi_1\circ f\circ\varphi_1^{-1}
 =(\psi_1\circ\psi^{-1})\circ(\psi\circ f\circ\varphi^{-1})\circ(\varphi\circ\varphi_1^{-1})
\]
是光滑的，因此 $f$ 仍为可微映射。

以上我们结束了有关微分流形及其上的映射的基本概念的讨论，下面就开始讨论微分流形上的微分学问题。
